\begin{topic}[id=quantum_intro,
             difficulty={Mittel},
             type={Überblick + Einordnung},
             prerequisites={Mathematische Grundkenntnisse; Bereitschaft zur konzeptionellen Einarbeitung},
             practice={optional},
             owner={Dozententeam},
             updated={2026-04-09},
             tags={ Quantencomputing, Qubits, Gate-Model, NISQ, Fehlerkorrektur }]{Quantencomputer – Stand \& Einordnung}

\TopicDescription{Quantencomputer versprechen Beschleunigung für bestimmte Problemklassen, sind aber in der NISQ-Ära noch stark begrenzt. Du lernst Grundbegriffe (Qubits, Gates, Messung), Hardwareansätze, Fehlermodelle und warum Fehlerkorrektur der zentrale Engpass ist. Ziel ist eine nüchterne Einordnung zwischen Hype und realem Nutzen – inkl. typischer Missverständnisse.}

\TopicLeitfragen{\item Welche Probleme profitieren wirklich?
\item Was limitiert aktuelle NISQ-Geräte?
\item Wie liest man Roadmaps kritisch?
}
\TopicScope{\item Keine Mathematik/Beweise.
\item Keine SDK-Tutorials.
}
\TopicPractice{Optionale Praxisidee: Kleines Beispiel (Qiskit/Cirq) zur Illustration von Rauscheffekten.}

\TopicKeywords{Quantencomputing, Qubits, Gate-Model, NISQ, Fehlerkorrektur}

\nocite{qiskitDocs,nistQcOverview}
\end{topic}
